% ============================================================
% CHAPTER 13 — Release: Publishing to the Stores
% ============================================================
\chapter{Release: Publishing to the Stores}\index{release}\index{Google Play}\index{App Store}

\section*{What You'll Build}
In this chapter, you will prepare the Flutter app for publication:
\begin{itemize}
  \item App icon and splash screen generated by AI
  \item Production build for Android (APK/AAB) and iOS (IPA)
  \item App digital signing
  \item Store listings (text, screenshots)
  \item Publishing process on Google Play and App Store
\end{itemize}

\textbf{Estimated time}: 90--120 minutes\\
\textbf{Prerequisite}: Working Flutter app (Ch.~11--12)

% ---------------------------------------------------------
\section{Preparing for Production}

\begin{praticabox}[title={Hands-on --- Update \texttt{\_CONTEXT.md} for release}]
Add a section to the context:
\begin{lstlisting}[language={}]
## Release

### App Identity
- Name: Notes App
- Android Bundle ID: com.example.notesapp
- iOS Bundle ID: com.example.notesapp
- Version: 1.0.0+1

### Environments
- Dev: http://10.0.2.2:3000/api (emulator)
- Staging: https://staging-api.notes-app.example.com/api
- Production: https://api.notes-app.example.com/api

### Build
- Android: flutter build appbundle (for Play Store)
- iOS: flutter build ipa (for App Store)
- Signing: keystore in android/app/keystore.jks (DO NOT commit)
\end{lstlisting}
\end{praticabox}

% ---------------------------------------------------------
\section{App Icon and Splash Screen}

\begin{praticabox}[title={Hands-on --- Generate the icon with AI}]
\begin{lstlisting}[language={}]
Generate an icon for the Notes app:
- Size: 1024x1024 pixels
- Background: indigo gradient (from #4F46E5 to #3730A3)
- Center icon: stylized note in white
- Style: flat, minimal, rounded corners
- Format: PNG with transparency

Save as assets/icon/app_icon.png
\end{lstlisting}
\end{praticabox}

\begin{tipbox}
If the AI cannot generate images directly, ask it to generate an SVG icon with code, or use tools like DALL-E or Midjourney with a prompt that the AI helps you write.
\end{tipbox}

\begin{praticabox}[title={Hands-on --- Configure flutter\_launcher\_icons}]
\begin{lstlisting}[language={}]
1. Add the dev dependency in pubspec.yaml
2. Add the configuration:
   flutter_launcher_icons:
     android: true
     ios: true
     image_path: "assets/icon/app_icon.png"
     adaptive_icon_background: "#4F46E5"
     adaptive_icon_foreground: "assets/icon/app_icon_foreground.png"
3. Generate the icons with: dart run flutter_launcher_icons
\end{lstlisting}
\end{praticabox}

\begin{praticabox}[title={Hands-on --- Splash Screen}]
\begin{lstlisting}[language={}]
1. Add flutter_native_splash in pubspec.yaml
2. Configure:
   flutter_native_splash:
     color: "#4F46E5"
     image: assets/icon/app_icon.png
     android_12:
       color: "#4F46E5"
       icon_background_color: "#4F46E5"
       image: assets/icon/app_icon.png
3. Generate with: dart run flutter_native_splash:create
\end{lstlisting}
\end{praticabox}

% ---------------------------------------------------------
\section{Production Configuration}

\begin{praticabox}[title={Hands-on --- Production environment}]
\begin{lstlisting}[language={}]
Modify api_config.dart to support environments:
1. Use --dart-define to pass the environment to the build
2. In production, the API URL points to the deployed backend
3. Disable debug logs in production
4. Configure: flutter build appbundle --dart-define=ENV=production
\end{lstlisting}
\end{praticabox}

The AI should generate something like:

\begin{lstlisting}[language=Java,title={ApiConfig with environments}]
class ApiConfig {
  static const String environment = String.fromEnvironment(
    'ENV',
    defaultValue: 'development',
  );

  static String get baseUrl {
    switch (environment) {
      case 'production':
        return 'https://api.notes-app.example.com/api';
      case 'staging':
        return 'https://staging-api.notes-app.example.com/api';
      default:
        return 'http://10.0.2.2:3000/api';
    }
  }
}
\end{lstlisting}

% ---------------------------------------------------------
\section{Android Build}

\subsection{Generating the Keystore}

The keystore is the digital certificate that signs your app. Without it, Google Play will not accept it.

\begin{lstlisting}[language=bash]
keytool -genkey -v -keystore android/app/keystore.jks \
  -keyalg RSA -keysize 2048 -validity 10000 \
  -alias notes-app
\end{lstlisting}

\begin{warnbox}
The \texttt{keystore.jks} file is \textsc{secret}. \textsc{Do not} commit it to git. If you lose it, you will never be able to update your app on the store again. Make a secure backup.
\end{warnbox}

\begin{praticabox}[title={Hands-on --- Configure signing}]
\begin{lstlisting}[language={}]
1. Create android/key.properties with the keystore data
   (storePassword, keyPassword, keyAlias, storeFile)
2. Modify android/app/build.gradle to read key.properties
3. Add key.properties and keystore.jks to .gitignore
4. Configure: minifyEnabled true, shrinkResources true
\end{lstlisting}
\end{praticabox}

\subsection{Build}

\begin{lstlisting}[language=bash]
# App Bundle (for Google Play --- recommended)
flutter build appbundle --dart-define=ENV=production

# APK (for direct distribution)
flutter build apk --release --dart-define=ENV=production
\end{lstlisting}

The generated file is located in:
\begin{itemize}
  \item \textbf{AAB}: \texttt{build/app/outputs/bundle/release/app-release.aab}
  \item \textbf{APK}: \texttt{build/app/outputs/flutter-apk/app-release.apk}
\end{itemize}

% ---------------------------------------------------------
\section{Build iOS}

\begin{warnbox}
To build for iOS, you need a Mac with Xcode installed and an Apple Developer account (\$99/year).
\end{warnbox}

\begin{praticabox}[title={Hands-on --- Xcode configuration}]
\begin{lstlisting}[language={}]
1. Open ios/Runner.xcworkspace in Xcode
2. In Signing & Capabilities:
   - Select your Team (Apple Developer account)
   - Verify the Bundle Identifier
3. In General:
   - Set the version (1.0.0)
   - Set the build number (1)
4. In Build Settings:
   - iOS Deployment Target: 15.0 (or higher)
\end{lstlisting}
\end{praticabox}

\begin{lstlisting}[language=bash]
flutter build ipa --dart-define=ENV=production
\end{lstlisting}

The IPA file is located in \texttt{build/ios/ipa/}.

% ---------------------------------------------------------
\section{Store Listings}

\begin{praticabox}[title={Hands-on --- Generate the text with AI}]
\begin{lstlisting}[language={}]
Generate the text for publishing on the stores:
1. Title: maximum 30 characters
2. Short description: maximum 80 characters
3. Full description: 4000 characters max
   - Introductory paragraph
   - Main features list
   - Security/privacy section
   - Final call to action
4. Keywords: 7-10 relevant keywords
5. Category: Productivity
6. Content rating: Everyone / 4+

Generate both the Italian and English versions.
\end{lstlisting}
\end{praticabox}

For screenshots, run the app on the emulator and capture the main screens:

\begin{enumerate}
  \item \textbf{Login screen} --- with the OAuth buttons
  \item \textbf{Dashboard} --- with some sample notes
  \item \textbf{Note detail} --- the open note
  \item \textbf{Create note} --- the entry form
  \item \textbf{Search} --- search results
\end{enumerate}

\begin{tipbox}
Use \texttt{flutter screenshot} to capture screenshots from the emulator directly from the terminal.
\end{tipbox}

% ---------------------------------------------------------
\section{Publishing on Google Play}

\begin{enumerate}
  \item Go to Google Play Console
  \item Create a new app
  \item Fill in the listing page (text, screenshots, icon)
  \item Go to Release $\rightarrow$ Production $\rightarrow$ Create new release
  \item Upload the \texttt{.aab} file
  \item Complete the content compliance declaration
  \item Set the price (free) and distribution countries
  \item Submit for review
\end{enumerate}

The first review generally takes 1--3 business days.

\subsection{Privacy Policy}

Google Play requires a privacy policy. Ask the AI:

\begin{praticabox}[title={Hands-on --- Privacy policy}]
\begin{lstlisting}[language={}]
Generate a privacy policy for the Notes app that:
- Explains that the app uses OAuth for authentication
- Note data is saved on secure servers
- Does not share data with third parties
- The user can delete the account and all data
- Compliant with GDPR
- URL: https://notes-app.example.com/privacy
\end{lstlisting}
\end{praticabox}

% ---------------------------------------------------------
\section{Publishing on App Store}

\begin{enumerate}
  \item Go to App Store Connect
  \item Create a new app
  \item Fill in the information (text, screenshots, icon)
  \item In Xcode: Product $\rightarrow$ Archive $\rightarrow$ Distribute App $\rightarrow$ App Store Connect
  \item Go back to App Store Connect, select the build
  \item Submit for review
\end{enumerate}

Apple review generally takes 1--2 business days.

\begin{warnbox}
Apple is more restrictive than Google. Make sure that:
\begin{itemize}
  \item The app works without crashes
  \item All links (privacy policy, terms) are reachable
  \item The app has value for the user (it is not just a ``demo'')
  \item OAuth login works correctly
\end{itemize}
\end{warnbox}

% ---------------------------------------------------------
\section{CI/CD for Mobile}

\begin{praticabox}[title={Hands-on --- GitHub Actions for automated builds}]
\begin{lstlisting}[language={}]
Create a GitHub Actions workflow that:
1. Triggers on push to the main branch
2. Builds the Flutter app for Android
3. Runs flutter analyze and flutter test
4. Generates the release APK
5. Uploads the APK as an artifact

File: .github/workflows/flutter-build.yml
\end{lstlisting}
\end{praticabox}

The AI should generate something like:

\begin{lstlisting}[language={},title={flutter-build.yml}]
name: Flutter Build

on:
  push:
    branches: [main]
  pull_request:
    branches: [main]

jobs:
  build:
    runs-on: ubuntu-latest
    steps:
      - uses: actions/checkout@v4
      - uses: subosito/flutter-action@v2
        with:
          flutter-version: '3.x'
          channel: 'stable'
      
      - name: Install dependencies
        working-directory: notes_mobile
        run: flutter pub get
      
      - name: Analyze
        working-directory: notes_mobile
        run: flutter analyze
      
      - name: Test
        working-directory: notes_mobile
        run: flutter test
      
      - name: Build APK
        working-directory: notes_mobile
        run: flutter build apk --release
             --dart-define=ENV=production
      
      - name: Upload APK
        uses: actions/upload-artifact@v4
        with:
          name: release-apk
          path: notes_mobile/build/app/outputs/
                flutter-apk/app-release.apk
\end{lstlisting}

\begin{notebox}
For automatic publication on the Play Store, you can add the Fastlane plugin to the workflow. But for the first release, manual publication is sufficient.
\end{notebox}

% ---------------------------------------------------------
\section*{What You Built --- Final Project}

In 13~chapters, you built a \textbf{complete system}:

\begin{figure}[H]
\centering
\includegraphics[width=0.85\textwidth]{06-notes-fullstack-arch}
\caption{Final Notes Full-Stack architecture}
\end{figure}

% ---------------------------------------------------------
\section*{Summary}

\begin{center}
\begin{tabularx}{\textwidth}{l X}
\toprule
\textbf{Aspect} & \textbf{Detail} \\
\midrule
Icon & flutter\_launcher\_icons with adaptive icon \\
Splash & flutter\_native\_splash with indigo theme \\
Android build & AAB signed with keystore \\
iOS build & IPA via Xcode (requires Mac + Apple Developer) \\
CI/CD & GitHub Actions: analyze + test + build \\
Store listing & Text generated by AI in IT and EN \\
Privacy & GDPR policy generated by AI \\
\bottomrule
\end{tabularx}
\end{center}
